\documentclass{crypto-exercise}
\usepackage{hyperref}
\author{Sven Laur}

\tags{Computational Diffie-Hellman problem, Decisional Diffie-Hellman problem, bilinear pairing}

\begin{document}
\begin{exercise}{Forced equivalence between CDH and DDH}
  Let $\GG$ be a finite multiplicative group of prime order $q$ such that all elements $y\in\GG$ can be
  expressed as powers of of $g\in\GG$.
  \begin{itemize}
  \item  Then the Computational
  Diffie-Hellman (CDH) problem is following. Given $u=g^x$ and
  $v=g^y$, find a group element $w=g^{xy}$. 
  \item Then the Decisional
  Diffie-Hellman (DDH) problem is the following. For any triple
  $u,v,w\in\GG$, you must decide whether it is a Diffie-Hellman triple
  or not.
  \item Let $\OOO$ be a sealed CDH oracle that operates as follows. Upon its first invocation, it issues $\ell$ DDH challenges. That is, for each challenge it samples a random tuple $(u,v,w)\gets \GG^3$ with probability $\tfrac{1}{2}$, and a valid Diffie--Hellman tuple otherwise. After each challenge, the oracle expects a reply bit $b$.
If all $\ell$ replies are correct, then $\OOO$ begins to behave as a CDH oracle --  it answers queries according to
\begin{align*}
\OOO(g^x, g^y) = g^{xy}.
\end{align*}
\end{itemize}
  Prove that if there exists an algorithm $\AD$ that has the advantage $\ADVDDH{\GG}{\AD}=\varepsilon$ then there exists an efficient oracle algorithm $\ADB^\OOO$ that achives significant $\ADVCDH{\GG}{\ADB^\OOO}$  even if the group $\GG$ was originally  a secure CDH group. This shows that we cannot ignore the possibility that the security parameters for the CDH and DDH problems may be of comparable magnitude. 
  \end{exercise}
\begin{solution}

\end{solution}
\end{document}
